% 申请人资料来自 docs/applicants/yan-feng.md；专业邮箱、办公地点已对照个人主页。
% 专业邮箱及通信地址用于当前草案，正式提交前由申请人终核。
% 性别、出生年月和联系电话按申请人共用资料及本轮补充授权填写；不填证件号码。
% 2026-09-10 由申请人本人提供：民族汉族、国别中国、每年工作时间 7 个月。传真仍留空。
% 2026-09-09 核对学院官网的华瑞街 777 号、邮编 830046 及教师页 B201 办公室。
% 地址按既有单行位点精简：省级与区级行政前缀省略，街道门牌、单位、房间和邮编保留。
% 来源：https://it.xju.edu.cn/info/1144/3131.htm；https://it.xju.edu.cn/xygk/xyjj.htm
% 中文工作题目沿用旧选题；英文题目保留研究含义并缩短以适合原有单行位点。
% 关键词按现有研究内容填写。
% 2026 开放基金官网指南及用户选择：一般课题，申请总额按 3.00 万元拟稿。
% 表格栏内只填金额；“预算草案”属写作阶段的内部说明，不随正式表格提交。
% 正文与配图已按 2026.11--2027.10 一年执行期编排，与本表期限一致。
\CasmSetup{
  title={不完备描述下地理实体关系约束的遥感图文检索方法研究},
  title-en={Relation-Constrained Remote Sensing Retrieval from Incomplete Text},
  applicant={颜丰}, organization={新疆大学},
  sex={男}, birth={1990.11}, ethnicity={汉族}, degree={博士}, professional-title={副教授},
  work-months={7}, phone={13122045162}, email={yf@xju.edu.cn}, fax={}, country={中国},
  address={乌鲁木齐市华瑞街777号，新疆大学计算机学院B201（830046）},
  research-area={多模态学习、跨模态检索与遥感视觉理解},
  direction={全域遥感与智能处理},
  period={2026 年 11 月 1 日—2027 年 10 月 31 日},
  category={一般课题}, funding={3.00 万元},
  keywords={遥感图文检索；地理实体关系；不完备描述；跨模态学习}
}
